\originalchapter{角动量、交换对称与壳层}\label{ch:angular}
\originalsection{角动量代数}
空间转动的生成元满足 $[L_i,L_j]=\ii\hbar\epsilon_{ijk}L_k$. 因为 $[L^2,L_i]=0$, 可以同时指定 $l$ 与一个分量 $m$: $L^2|l,m\rangle=\hbar^2l(l+1)|l,m\rangle$, $L_z|l,m\rangle=\hbar m|l,m\rangle$. 三个分量不对易, 所以不能同时具有确定值.

令 $L_\pm=L_x\pm\ii L_y$, 则 $[L_z,L_\pm]=\pm\hbar L_\pm$. 这说明 $L_\pm$ 把 $m$ 变成 $m\pm1$. 用 $L_-L_+=L^2-L_z^2-\hbar L_z$ 计算作用后的范数, 得
\begin{equation}
 L_\pm|l,m\rangle=\hbar\sqrt{l(l+1)-m(m\pm1)}\,|l,m\pm1\rangle.
\end{equation}
有限维转动表示在最高、最低权处终止, 给 $m=-l,-l+1,\ldots,l$. 对轨道角动量, 单值球面函数要求 $l$ 为非负整数; 自旋也允许半整数, 不来自小球实际自转.

\originalsection{两个角动量的耦合}
当 $\mathbf J=\mathbf L+\mathbf S$ 时, 不耦合基为 $|l,m_l\rangle|s,m_s\rangle$, 耦合基为 $|j,m_j;l,s\rangle$. 二者只是同一态空间的不同正交基, Clebsch--Gordan 系数给出基变换. $m_j=m_l+m_s$, $|l-s|\le j\le l+s$, 相邻 $j$ 相差一.

两个自旋 $1/2$ 的四维空间分解为三重态和单重态:
\begin{align*}
 |1,1\rangle&=|\uparrow\uparrow\rangle,\quad |1,0\rangle=(|\uparrow\downarrow\rangle+|\downarrow\uparrow\rangle)/\sqrt2,\\
 |1,-1\rangle&=|\downarrow\downarrow\rangle,\quad |0,0\rangle=(|\uparrow\downarrow\rangle-|\downarrow\uparrow\rangle)/\sqrt2.
\end{align*}
前者在交换粒子时对称, 后者反对称. 用 $\mathbf L\cdot\mathbf S=(J^2-L^2-S^2)/2$, $s=1/2$ 的自旋轨道相互作用 $V_{\rm so}\mathbf L\cdot\mathbf S/\hbar^2$ 给出
\begin{equation}
 E_{l+1/2}=E_0+\frac l2V_{\rm so},\qquad E_{l-1/2}=E_0-\frac{l+1}{2}V_{\rm so}.
\end{equation}
核中通常相应系数使 $j=l+1/2$ 更低, 与简单三维振子不同. 加权重心不变, 因为简并度 $2l+2$ 与 $2l$ 乘能移后相消.

\originalsection{费米子与核壳模型}
同类核子的总态在交换时必须反对称. 两个中子的同位旋态对称, 若低能轨道为 $l=0$ 则空间态也对称, 因此自旋只能是反对称的单重态. 这排除把氘核的三重态相互作用直接搬到双中子体系; 是否束缚还须依赖单重态相互作用的实验强度, 不能由 Pauli 原理独自证明.

在平均场近似下, 核子占据单粒子轨道 $nljm_j$, 每轨道容纳 $2j+1$ 个同类核子. 球对称振子若无自旋轨道项, 按主振子数 $N=2n_r+l$ 分类, 每种核子的壳容量为 $(N+1)(N+2)$. 得到 $2,8,20,40,70,\ldots$, 不能解释全部实验魔数 $2,8,20,28,50,82,126$. 强自旋轨道分裂将大 $j$ 轨道下移, 改变壳闭合次序.

偶偶核倾向基态 $J^\pi=0^+$; 奇 $A$ 核在简单价核子模型中常由未成对粒子或空穴确定 $j$ 和宇称 $(-1)^l$. ``空穴''指闭壳少一个粒子, 它具有相同 $j$ 和宇称标记, 但不是把能量直接取负. 配对、残余相互作用和形变可改变轨道次序, 因此价核子模型对激发态尤其要慎用.

\originalsection{原子模型与核模型的关系}
氢样原子中 $V=-Ze^2/(4\pi\epsilon_0r)$ 产生 $E_n=-\mu Z^2e^4/[2(4\pi\epsilon_0)^2\hbar^2n^2]$. 核平均场则短程且在内部较平坦, 常用有限深球阱或 Woods--Saxon 势. 二者同用球谐函数分离角部分, 径向谱却完全不同. 核物理不能以原子 Bohr 半径替代 fm 尺度.

三维振子能量 $E_N=\hbar\omega(N+3/2)$ 来自三个独立一维振子; 核经验近似常取 $\hbar\omega\simeq41A^{-1/3}\,\mathrm{MeV}$. 这是平均场标度, 不是任意核的精确激发谱. 从单粒子能量直接求总结合能还存在相互作用双计数问题: 每一对粒子的势能在两个单粒子本征值中出现两次. 作业中的简单求和应明确它是指定的模型估算.
